
\documentclass{article}
\usepackage{colortbl}
\usepackage{makecell}
\usepackage{multirow}
\usepackage{supertabular}

\begin{document}

\newcounter{utterance}

\centering \large Interaction Transcript for game `ifeval', experiment `ifeval', episode 154 with refineonly{-}llama{-}3.1{-}8b{-}instruct.
\vspace{24pt}

{ \footnotesize  \setcounter{utterance}{1}
\setlength{\tabcolsep}{0pt}
\begin{supertabular}{c@{$\;$}|p{.15\linewidth}@{}p{.15\linewidth}p{.15\linewidth}p{.15\linewidth}p{.15\linewidth}p{.15\linewidth}}
    \# & \multicolumn{2}{c}{Player} && \multicolumn{2}{c}{Game Master} \\
    \hline

    \theutterance \stepcounter{utterance}  
    & & & \multicolumn{4}{p{0.6\linewidth}}{
        \cellcolor[rgb]{0.9,0.9,0.9}{
            \makecell[{{p{\linewidth}}}]{
                \texttt{\tiny{[P1$\langle$GM]}}
                \texttt{Explain Generative Adversarial Networks (GANs) to me using bullet points. Do not contain any commas in your response. End your response with a postscript indicated by P.P.S} \\
\texttt{Include the keywords "lacking", "model", "performance", "quality", "architecture".} \\
            }
        }
    }
     \\ \\

    \theutterance \stepcounter{utterance}  
    & & \multicolumn{4}{p{0.6\linewidth}}{
        \cellcolor[rgb]{0.95,0.95,0.95}{
            \makecell[{{p{\linewidth}}}]{
                \texttt{\tiny{[P1$\leftrightarrow$P]}}
                \texttt{{[}DRAFT{]}} \\
\texttt{Here's an explanation of Generative Adversarial Networks (GANs) using bullet points:} \\
\\ 
\texttt{* GANs are a type of deep learning model that consists of two neural networks: a generator and a discriminator.} \\
\texttt{* The generator model takes a random noise vector as input and produces a synthetic data sample that is similar to real data.} \\
\texttt{* The discriminator model takes a data sample as input and outputs a probability that the sample is real.} \\
\texttt{* The generator and discriminator models are trained simultaneously in a competitive process.} \\
\texttt{* The generator tries to produce synthetic data samples that are indistinguishable from real data samples.} \\
\texttt{* The discriminator tries to correctly classify real and synthetic data samples.} \\
\texttt{* The performance of the generator and discriminator models improves as they are trained together.} \\
\texttt{* The quality of the generated data samples improves as the generator model becomes more sophisticated.} \\
\texttt{* The architecture of the generator and discriminator models can be varied to improve performance.} \\
\texttt{* GANs are lacking in terms of interpretability and stability but have shown impressive results in image and video generation.} \\
\texttt{* GANs have many applications including data augmentation image and video editing and style transfer.} \\
\texttt{* The training process of GANs can be challenging and requires careful tuning of hyperparameters.} \\
\\ 
\texttt{P.P.S} \\
            }
        }
    }
    & \\ \\

    \theutterance \stepcounter{utterance}  
    & \multicolumn{4}{p{0.6\linewidth}}{
        \cellcolor[rgb]{0.561,0.737,0.561}{
            \makecell[{{p{\linewidth}}}]{
                \texttt{\tiny{[P1$\rangle$GM]}}
                \texttt{* GANs are a type of deep learning model that consists of two neural networks a generator and a discriminator.} \\
\texttt{* The generator model takes a random noise vector as input and produces a synthetic data sample that is similar to real data.} \\
\texttt{* The discriminator model takes a data sample as input and outputs a probability that the sample is real.} \\
\texttt{* The generator and discriminator models are trained simultaneously in a competitive process.} \\
\texttt{* The generator tries to produce synthetic data samples that are indistinguishable from real data samples.} \\
\texttt{* The discriminator tries to correctly classify real and synthetic data samples.} \\
\texttt{* The performance of the generator and discriminator models improves as they are trained together.} \\
\texttt{* The quality of the generated data samples improves as the generator model becomes more sophisticated.} \\
\texttt{* The architecture of the generator and discriminator models can be varied to improve performance.} \\
\texttt{* GANs are lacking in terms of interpretability and stability but have shown impressive results in image and video generation.} \\
\texttt{* GANs have many applications including data augmentation image and video editing and style transfer.} \\
\texttt{* The training process of GANs can be challenging and requires careful tuning of hyperparameters.} \\
\texttt{* A well{-}designed model is crucial for achieving good performance and quality in GANs.} \\
\texttt{* The generator model needs to be able to produce high{-}quality data samples to fool the discriminator.} \\
\texttt{* The discriminator model needs to be able to correctly classify real and synthetic data samples to improve the performance of the generator.} \\
\texttt{* The quality of the generated data samples is directly related to the performance of the generator model.} \\
\texttt{* A good architecture for the generator and discriminator models is essential for achieving good performance and quality in GANs.} \\
\\ 
\texttt{P.P.S} \\
            }
        }
    }
    & & \\ \\

    \theutterance \stepcounter{utterance}  
    & & & \multicolumn{2}{p{0.3\linewidth}}{
        \cellcolor[rgb]{0.9,0.9,0.9}{
            \makecell[{{p{\linewidth}}}]{
                \texttt{\tiny{[GM$|$GM]}}
                {'punctuation:no_comma': {}, 'detectable_content:postscript': {'postscript_marker': 'P.P.S'}, 'keywords:existence': {'keywords': ['lacking', 'model', 'performance', 'quality', 'architecture']}}
            }
        }
    }
    & & \\ \\

    \theutterance \stepcounter{utterance}  
    & & & \multicolumn{2}{p{0.3\linewidth}}{
        \cellcolor[rgb]{0.9,0.9,0.9}{
            \makecell[{{p{\linewidth}}}]{
                \texttt{\tiny{[GM$|$GM]}}
                \texttt{game\_result = WIN} \\
            }
        }
    }
    & & \\ \\

\end{supertabular}
}

\end{document}
